\section{Certificate and reproducibility boundary}

The finite part of Theorem~\ref{thm:nonsurjective} is deliberately small enough to be reproduced without the larger compressed census used elsewhere in the project. All files named below are in the source repository \cite{repository}, which also pins the toolchain versions used to produce them.

\subsection*{Certificate}
The committed file
\[
 \texttt{certificates/spectrum\_m259.csv}
\]
contains 259 rows, one for each $1\le m<260$, with columns
\[
 (m,t,c,b_t).
\]
Its SHA-256 digest is
\begin{center}
\small\texttt{66a06cff15735c4a3caf98575f29afbcd881fbef06334616fbc3bc772b7ab084}.
\end{center}

\subsection*{Independent regeneration}
The script
\[
 \texttt{independent/verify\_small\_spectrum.py}
\]
uses arbitrary-precision integers and the literal recurrence \eqref{eq:recurrence}. For each start it:
\begin{enumerate}
\item iterates until the first index satisfying the absorption criterion;
\item records $(m,t,c,b_t)$;
\item confirms 64 subsequent increments equal $c$;
\item regenerates the complete CSV and checks its digest;
\item asserts that $5$ and $7$ are absent and that the witnesses in Table~\ref{tab:witnesses} occur.
\end{enumerate}

The reproduction commands are
\begin{verbatim}
python independent/verify_small_spectrum.py
sha256sum certificates/spectrum_m259.csv
\end{verbatim}
On Windows PowerShell, the second command may be replaced by
\begin{verbatim}
Get-FileHash -Algorithm SHA256 certificates\spectrum_m259.csv
\end{verbatim}

\subsection*{Logical boundary}
The proof of the finite-start bound is symbolic mathematics. The certificate supplies only the finite exhaustion authorized by that theorem. The nonsurjectivity conclusion does not rely on the project's $10^7$-start compressed census, on any heuristic model, or on universal stabilization.

The Lean development in the repository formalizes the foundational identities of Sections 2 and 3, both load-bearing steps of Section 4 — the forced-rebound lemma and the ratchet, including the induction of \eqref{eq:ratchet-invariant} — and Theorem~\ref{thm:finite-start} itself, as \texttt{Conjecture.finite\_start}. The entry lemma is formalized in the form the proof consumes, namely existence of an index with $b_n<n^2$ together with minimality; the explicit bound $n_0\le\lceil\sqrt{2m}\,\rceil+2$ is not formalized and is not needed, since the proof of Theorem~\ref{thm:finite-start} never uses it. The quadratic step $(n_0-1)^2<(c+2)n_0\Rightarrow n_0\le c+4$ is formalized as \texttt{Conjecture.quadratic\_step}.

Two boundaries should be stated precisely. First, the formalized theorem takes as hypothesis that the orbit \emph{lies on the absorbing ray} at some index, $b_t=c(t+1)$ with $c<t$. By the absorption criterion this is equivalent to having eventual increment $c$, but only the direction ray $\Rightarrow$ constant increment is formalized; the converse uses the limiting divisibility argument in the proof of that criterion and is not. So the machine-checked statement is Theorem~\ref{thm:finite-start} from the ray hypothesis, and a fully formal route from ``the orbit stabilizes with increment $c$'' remains open. Second, Theorem~\ref{thm:nonsurjective} is not formalized: its finite exhaustion is checked by an independent arbitrary-precision program rather than inside the proof assistant.
